\begin{table*}[t]
\centering
\small
\setlength{\tabcolsep}{4pt}
\renewcommand{\arraystretch}{1.08}
\resizebox{\textwidth}{!}{%
\begin{tabular}{p{0.17\linewidth}p{0.25\linewidth}p{0.12\linewidth}p{0.12\linewidth}p{0.26\linewidth}}
\toprule
Case & Prompt signal & Gold action & Model action & Action-selection lesson \\
\midrule
Stale premise & User asks why Facebook is still trading under \texttt{FB}; evidence says Meta changed the ticker to \texttt{META} on June 9, 2022. & \texttt{challenge} & \texttt{answer} & The model mentions the correction but still answers through the stale premise. Factual knowledge is present, but the next action is wrong. \\
False premise & A math-solution step uses $1 + 2/(3n+1)$ where the supported step is $1 + 1/(3n+1)$. & \texttt{challenge} & \texttt{answer} & The model explains why logarithms are useful instead of interrupting the defective step. Fluent explanation can rationalize a bad premise. \\
Answerable control & Jenny has 36 cookies; her son eats 20 and her husband eats 7. & \texttt{answer} & \texttt{challenge} & The supported answer is 9, but the model claims the prompt is incomplete. The task does not reward caution when the input is already answerable. \\
Conflict but answerable & Noisy passages ask who was told to beware the Ides of March. Several passages still support Caesar. & \texttt{answer} & \texttt{answer} & Retrieval noise alone is not an abstention condition. The action depends on whether enough support remains. \\
\bottomrule
\end{tabular}
}
\caption{Representative action-boundary examples from the Day-1 evidence bundle. The table reports the next-action label rather than grading only the final answer text.}
\label{tab:qualitative-examples}
\end{table*}


Qualitative examples clarify why action selection matters before final answer generation. Table~\ref{tab:qualitative-examples} shows four boundary cases from the active evidence bundle. In a stale-premise example, the user asks why Facebook is still trading under \texttt{FB}; the correct next action is \texttt{challenge}, because Meta's stock began trading under \texttt{META} in 2022. Qwen2.5-1.5B instead answers while mentioning the correction. The utility loss is not lack of factual knowledge; it is accepting the user's outdated framing as an answer request \citep{kasai-etal-2023-realtimeqa,vu-etal-2023-freshllms}.

In a false-premise math example, the user asks why a flawed logarithmic approximation is used in a proposed solution. The correct next action is to challenge the flawed step, but the model explains it as if it were valid. This is the central action-boundary failure for reasoning-style responses: fluent intermediate reasoning can rationalize a bad premise instead of interrupting it \citep{lin-etal-2022-truthfulqa,ji-etal-2023-hallucination-survey}.

The 7B reasoning checkpoint shows the complementary risk. On a clean answerable triangulation problem, it produces extended reasoning but predicts \texttt{abstain} rather than committing to the supported answer. On a false-premise tournament problem, it reasons about the stated premise but does not explicitly correct the false $3^t$ construction to the supported $2^t$ construction. These cases demonstrate that the benchmark does not reward caution or reasoning length alone; the model must choose the right repair action for the input state \citep{kamath-etal-2020-selective,kadavath-etal-2022-language}. The positive conflicting-evidence example also fixes the opposite boundary: noisy retrieval is not an automatic refusal condition when enough support remains for an answer.
